% Reviewed translation: PDF pages 66--70 / printed pages 52--56.
% Stops on PDF page 70 before Section 4 begins.

\MathBookSourcePage{66}
\setcounter{statement}{3}
\begin{Lemma}\label{lem:sec3-w0-curl-isomorphism}
这里只假设 $\Omega$ 有界、Lipschitz 连续、单连通, 且 $\Gamma$ 只有一个连通分支
($p=0$). 则映射 $\boldsymbol\phi\mapsto\Curl\boldsymbol\phi$ 是从 $W_0$ 到 $H$
的同构, 并且存在两个正常数 $C_1,C_2$, 使得
\begin{equation}\label{eq:sec3-w0-curl-poincare}
\|\boldsymbol\phi\|_{0,\Omega}
\le C_1\|\Curl\boldsymbol\phi\|_{0,\Omega}
\qquad\forall\boldsymbol\phi\in W_0,
\end{equation}
\begin{equation}\label{eq:sec3-w-div-curl-poincare}
\|\boldsymbol\phi\|_{0,\Omega}
\le C_2\{\|\Curl\boldsymbol\phi\|_{0,\Omega}
+\|\Div\boldsymbol\phi\|_{0,\Omega}\}
\qquad\forall\boldsymbol\phi\in W.
\end{equation}
\end{Lemma}
\begin{Proof}
因为 $\Omega$ 单连通且 $p=0$, Theorem~\ref{thm:sec3-tangential-zero-vector-potential}
断言 $H$ 中每个函数 $\mathbf u$ 在 $W_0$ 中恰有一个向量势 $\boldsymbol\phi$:
\[
\mathbf u=\Curl\boldsymbol\phi.
\]
因此映射 $\Curl$ 是从 $W_0$ 到 $H$ 的一一、线性、连续的满射. 由于 $W_0$ 和 $H$
都是 Banach 空间, 它是同构. 此外, 映射
\[
\boldsymbol\phi\mapsto\|\mathbf u\|_{H(\Div;\Omega)}
=\|\Curl\boldsymbol\phi\|_{0,\Omega}
\]
是 $W_0$ 上与 $W$ 的范数等价的范数. 这就证明了
\eqref{eq:sec3-w0-curl-poincare}.

现在, \eqref{eq:sec3-w-div-curl-poincare} 是
\eqref{eq:sec3-w0-curl-poincare} 与 Lemma~\ref{lem:sec3-w-embedding-reduction}
中分解的简单推论: $W$ 中的 $\boldsymbol\phi$ 可写为
\[
\boldsymbol\phi=\boldsymbol\psi+\Grad q,
\]
其中 $\boldsymbol\psi\in W_0$, 且由 Poincaré Theorem~\ref{thm:poincare-friedrichs},
\[
|q|_{1,\Omega}\le C\|\Div\boldsymbol\phi\|_{0,\Omega}.
\]
所以
\[
\|\boldsymbol\phi\|_{0,\Omega}
\le C_1\|\Curl\boldsymbol\psi\|_{0,\Omega}
+C\|\Div\boldsymbol\phi\|_{0,\Omega},
\]
而 $\Curl\boldsymbol\psi=\Curl\boldsymbol\phi$, 结论随即得到.
注意, 虽然这里使用 Lemma~\ref{lem:sec3-w-embedding-reduction} 证明的一部分,
但并不需要其中关于 $\Omega$ 的假设, 因为这里不要求 $q\in H^2(\Omega)$.
\end{Proof}

\setcounter{statement}{6}
\begin{Theorem}\label{thm:sec3-w-h1-embedding}
空间 $W$ 连续嵌入 $H^1(\Omega)^3$.
\end{Theorem}
\begin{Proof}
$1^\circ)$ 暂且假设 $\Omega$ 单连通且 $p=0$. 根据
Lemma~\ref{lem:sec3-w-embedding-reduction}, 只需考察 $W_0$.
由 Lemma~\ref{lem:sec3-w0-curl-isomorphism}, 映射 $\Curl$ 是从 $W_0$ 到 $H$ 的同构.
而 Theorem~\ref{thm:sec3-tangential-zero-vector-potential} 的假设满足, 所以
$\Curl$ 也是从 $W_0\cap H^1(\Omega)^3$ 到 $H$ 的同构. 因而在代数意义和拓扑意义下
\[
W_0=W_0\cap H^1(\Omega)^3,
\]
所以 $W_0$ 连续嵌入 $H^1(\Omega)^3$.

$2^\circ)$ 当 $\Omega$ 任意时, 它不再是凸的, 因而由假设 $\Gamma$ 必须属于
$\mathcal C^{1,1}$ 类. 此时容易找到 $\Omega$ 的有限开覆盖
$\{\mathcal O_i\}_{1\le i\le l}$, 使得每个 $\mathcal O_i\cap\Omega$ 单连通且边界连通.
令 $\{\alpha_i\}$ 是从属于 $\{\mathcal O_i\}$ 的单位分解; 那么 $W$ 中每个
$\boldsymbol\phi$ 都可表示为

\MathBookSourcePage{67}
\[
\boldsymbol\phi=\sum_{i=1}^l\alpha_i\boldsymbol\phi
\quad\text{in }\Omega.
\]
因此问题归结为 $\alpha_i\boldsymbol\phi$ 的正则性. 注意, 因为 $\Gamma$ 光滑且
$\alpha_i\in\D(\mathcal O_i)$, 可以认为 $\alpha_i\boldsymbol\phi$ 定义在
$\mathcal O_i\cap\Omega$ 的一个子集 $\mathcal O_i'$ 上, 该子集具有光滑边界,
并具有与 $\mathcal O_i\cap\Omega$ 相同的性质. 此外,
\[
\alpha_i(\boldsymbol\phi\times\mathbf n)=\mathbf0
\quad\text{on }\partial\mathcal O_i',
\]
因为在 $\partial\mathcal O_i'$ 的每一点上至少有一个因子为零.
所以 $\alpha_i\boldsymbol\phi\in W$ on $\mathcal O_i'$; 由 $1^\circ)$ 得
$\alpha_i\boldsymbol\phi\in H^1(\mathcal O_i')^3$, 且
\[
\|\alpha_i\boldsymbol\phi\|_1\le C_i
\{\|\alpha_i\boldsymbol\phi\|_0
+\|\Div(\alpha_i\boldsymbol\phi)\|_0
+\|\Curl(\alpha_i\boldsymbol\phi)\|_0\},
\]
其中所有范数都取在 $\mathcal O_i'$ 上. 因此在相同记号下,
\[
\|\alpha_i\boldsymbol\phi\|_1\le C_i'
\{\|\boldsymbol\phi\|_0+\|\Div\boldsymbol\phi\|_0
+\|\Curl\boldsymbol\phi\|_0\},
\]
从而证明该 Theorem.
\end{Proof}

\setcounter{subsection}{4}
\subsection{$H_0(\Div;\Omega)\cap H(\Curl;\Omega)$ 到 $H^1(\Omega)^3$ 的嵌入}
\label{subsec:sec3-h0div-hcurl-embedding}

继续保持 Subsection~\ref{subsec:sec3-hdiv-h0curl-embedding} 的假设, 并置
\[
U=H_0(\Div;\Omega)\cap H(\Curl;\Omega),
\qquad
U_0=\{\mathbf u\in U;\ \Div\mathbf u=0\}.
\]
先给出与 Lemma~\ref{lem:sec3-w-embedding-reduction} 和
Lemma~\ref{lem:sec3-w0-curl-isomorphism} 类似的两个结果.

\setcounter{statement}{4}
\begin{Lemma}\label{lem:sec3-u-embedding-reduction}
空间 $U$ 连续嵌入 $H^1(\Omega)^3$, 当且仅当 $U_0$ 也连续嵌入
$H^1(\Omega)^3$.
\end{Lemma}

证明沿用 Lemma~\ref{lem:sec3-w-embedding-reduction} 的思路, 留给读者.

\setcounter{statement}{5}
\begin{Lemma}\label{lem:sec3-u0-curl-isomorphism}
设 $\Omega$ 有界、Lipschitz 连续且单连通. 则映射
$\boldsymbol\phi\mapsto\Curl\boldsymbol\phi$ 是从 $U_0$ 到空间
\[
T=\{\mathbf u\in L^2(\Omega)^3;\ \Div\mathbf u=0,
\ \langle\mathbf u\cdot\mathbf n,1\rangle_{\Gamma_i}=0,
\ 0\le i\le p\}
\]
的同构, 并且存在两个正常数 $C_1,C_2$, 使得
\begin{equation}\label{eq:sec3-u0-curl-poincare}
\|\boldsymbol\phi\|_{0,\Omega}
\le C_1\|\Curl\boldsymbol\phi\|_{0,\Omega}
\qquad\forall\boldsymbol\phi\in U_0,
\end{equation}
\begin{equation}\label{eq:sec3-u-div-curl-poincare}
\|\boldsymbol\phi\|_{0,\Omega}
\le C_2\{\|\Curl\boldsymbol\phi\|_{0,\Omega}
+\|\Div\boldsymbol\phi\|_{0,\Omega}\}
\qquad\forall\boldsymbol\phi\in U.
\end{equation}
\end{Lemma}
\begin{Proof}
因为 $\Omega$ 单连通, 由 Theorem~\ref{thm:sec3-normal-zero-vector-potential}
和 Lemma~\ref{lem:sec3-w0-curl-isomorphism} 的论证可知, 映射
$\boldsymbol\phi\mapsto\Curl\boldsymbol\phi$ 是从 $U_0$ 到 $T$ 的同构,
并且 \eqref{eq:sec3-u0-curl-poincare} 成立. 进而,
\eqref{eq:sec3-u-div-curl-poincare} 由
\eqref{eq:sec3-u0-curl-poincare}, 分解
\eqref{eq:sec3-l2-three-dimensional-decomposition} 和
Theorem~\ref{thm:h1-quotient-hilbert} 立即得到.
注意, 这里同样不需要 $\Gamma$ 有超出 Lipschitz 连续性的连续性条件.
\end{Proof}

\MathBookSourcePage{68}
\setcounter{statement}{6}
\begin{Lemma}\label{lem:sec3-u-embedding-neumann-criterion}
如果 $U$ 连续嵌入 $H^1(\Omega)^3$, 则对每个满足
$\int_\Gamma\boldsymbol\psi\cdot\mathbf n\dd s=0$ 的
$\boldsymbol\psi\in H^1(\Omega)^3$, 非齐次 Neumann 问题
\begin{equation}\label{eq:sec3-u-embedding-neumann-problem}
\Delta\chi=0\quad\text{in }\Omega,
\qquad
\partial\chi/\partial n=\boldsymbol\psi\cdot\mathbf n
\quad\text{on }\Gamma
\end{equation}
的解 $\chi$ 属于 $H^2(\Omega)/\mathbb R$. 反之, 当 $\Omega$ 单连通时,
这个条件足以保证上述嵌入.
\end{Lemma}
\begin{Proof}
$1^\circ)$ 设 $\boldsymbol\psi\in H^1(\Omega)^3$, 且
$\int_\Gamma\boldsymbol\psi\cdot\mathbf n\dd s=0$; 则
\eqref{eq:sec3-u-embedding-neumann-problem} 在 $H^1(\Omega)/\mathbb R$
中有唯一解 $\chi$, 并且
\[
\boldsymbol\phi=\boldsymbol\psi-\Grad\chi
\]
属于 $U$. 因此, 如果 $\boldsymbol\phi\in H^1(\Omega)^3$, 则必有
$\chi\in H^2(\Omega)/\mathbb R$.

$2^\circ)$ 反之, 由 Theorem~\ref{thm:sec3-normal-zero-vector-potential} 的证明可知,
如果 $\chi\in H^2(\Omega)/\mathbb R$, 则 $T$ 中每个函数在
$U_0\cap H^1(\Omega)^3$ 中都有唯一向量势. 因此 Lemma~\ref{lem:sec3-u-embedding-reduction}
和 Lemma~\ref{lem:sec3-u0-curl-isomorphism} 蕴含 $U$ 连续嵌入 $H^1(\Omega)^3$.
\end{Proof}

当 $\Gamma$ 属于 $\mathcal C^{1,1}$ 类时,
\eqref{eq:sec3-u-embedding-neumann-problem} 的解确实属于 $H^2(\Omega)/\mathbb R$;
所以当 $\Omega$ 单连通时, Lemma~\ref{lem:sec3-u-embedding-neumann-criterion}
保证 $U\subset H^1(\Omega)^3$. 此外, 使用
Theorem~\ref{thm:sec3-w-h1-embedding} 中 $2^\circ)$ 的论证可以去掉最后这个限制.
由此得到如下结果.

\setcounter{statement}{7}
\begin{Theorem}\label{thm:sec3-u-h1-smooth-boundary}
如果 $\Omega$ 有界且边界属于 $\mathcal C^{1,1}$ 类, 则空间 $U$ 连续嵌入
$H^1(\Omega)^3$.
\end{Theorem}

遗憾的是, 当 $\Gamma$ 是多面体且
$\boldsymbol\psi\in H^1(\Omega)^3$ 时, 它的法向分量
$\boldsymbol\psi\cdot\mathbf n$ 不属于 $H^{1/2}(\Gamma)$; 我们未能在文献中找到关于
\eqref{eq:sec3-u-embedding-neumann-problem} 的解的正则性的任何直接信息.
在这种情形下, Lemma~\ref{lem:sec3-u-embedding-neumann-criterion} 似乎无济于事;
我们提出用 Nédélec [60] 的另一种论证来建立 $U$ 到 $H^1(\Omega)^3$ 的嵌入.
简言之, 该作者使用一种特殊的范数等价关系, 它在光滑凸区域 $\Omega$ 上成立,
且其中的常数与 $\Omega$ 无关. 然后通过正则化凸多面体的边界, 把这种等价关系推广到凸多面体.
这种方法的出发点是下面的技术 Lemma; 证明太长, 无法在此收入, 可见 Grisvard [42].

\setcounter{statement}{7}
\begin{Lemma}\label{lem:sec3-curvature-div-curl-identity}
设 $\Omega$ 是 $\mathbb R^3$ 中的有界开子集, 边界 $\Gamma$ 属于
$\mathcal C^2$ 类. 每个满足 $\boldsymbol\phi\cdot\mathbf n=0$ on $\Gamma$ 的
$\boldsymbol\phi\in H^1(\Omega)^3$ 都有
\begin{equation}\label{eq:sec3-curvature-div-curl-identity}
|\boldsymbol\phi|_{1,\Omega}^2
+\int_\Gamma(\mathscr R\boldsymbol\phi,\boldsymbol\phi)\dd s
=\|\Curl\boldsymbol\phi\|_{0,\Omega}^2
+\|\Div\boldsymbol\phi\|_{0,\Omega}^2,
\end{equation}
其中 $\mathscr R$ 表示 $\Gamma$ 的切平面上的曲率张量.
\end{Lemma}

当 $\Omega$ 为凸集时, 曲率张量 $\mathscr R$ 正定, 因而立即得到如下范数等价.

\MathBookSourcePage{69}
\setcounter{statement}{5}
\begin{Corollary}\label{cor:sec3-convex-smooth-div-curl-bound}
设 $\Omega$ 是 $\mathbb R^3$ 中有界凸开子集, 边界 $\Gamma$ 属于
$\mathcal C^2$ 类. $U\cap H^1(\Omega)^3$ 中的函数 $\boldsymbol\phi$ 满足
\begin{equation}\label{eq:sec3-convex-smooth-div-curl-bound}
|\boldsymbol\phi|_{1,\Omega}^2
\le\|\Curl\boldsymbol\phi\|_{0,\Omega}^2
+\|\Div\boldsymbol\phi\|_{0,\Omega}^2.
\end{equation}
\end{Corollary}

\setcounter{statement}{8}
\begin{Theorem}\label{thm:sec3-convex-polyhedron-u-h1}
设 $\Omega$ 是有界凸多面体. 空间 $U$ 连续嵌入 $H^1(\Omega)^3$,
且 $U$ 中的函数满足 \eqref{eq:sec3-convex-smooth-div-curl-bound}.
\end{Theorem}
\begin{Proof}
可以构造(参见 Grisvard [42])一个递增的开凸集序列
$(\Omega_l)_{l\ge1}$, 每个集合都有 $\mathcal C^2$ 边界 $\Gamma_l$, 且
\[
\bar\Omega_l\subset\Omega\quad\forall l\ge1,
\qquad
\Omega=\bigcup_{l\ge1}\Omega_l.
\]
由 Lemma~\ref{lem:sec3-u-embedding-reduction}, 可以只考察 $U_0$.
因此, 令 $\boldsymbol\phi\in U_0$, 并令
$\chi_l\in H^1(\Omega_l)/\mathbb R$ 为问题
\begin{equation}\label{eq:sec3-inner-domain-neumann-problem}
\Delta\chi_l=0\quad\text{in }\Omega_l,
\qquad
\partial\chi_l/\partial n=\boldsymbol\phi\cdot\mathbf n
\quad\text{on }\Gamma_l
\end{equation}
的解; 在 $\Omega_l$ 上定义
\[
\boldsymbol\phi_l=\boldsymbol\phi-\Grad\chi_l.
\]
因为 $\Div\boldsymbol\phi=0$ in $\Omega$, 非齐次 Neumann 问题
\eqref{eq:sec3-inner-domain-neumann-problem} 有唯一解 $\chi_l$.
此外, 由 Corollary~\ref{cor:sec2-local-h1-hdiv-hcurl},
$\boldsymbol\phi\in H^1(\Omega_l)^3$ for all $l$; 又因为 $\Gamma_l$ 光滑,
这进一步蕴含 $\boldsymbol\phi\cdot\mathbf n\in H^{1/2}(\Gamma_l)$.
所以 $\chi_l\in H^2(\Omega_l)/\mathbb R$, 从而
$\boldsymbol\phi_l\in H^1(\Omega_l)^3$, 且
$\boldsymbol\phi_l\cdot\mathbf n=0$ on $\Gamma_l$.
因此可以把 Corollary~\ref{cor:sec3-convex-smooth-div-curl-bound} 应用于
$\boldsymbol\phi_l$:
\begin{equation}\label{eq:sec3-inner-domain-h1-bound}
|\boldsymbol\phi_l|_{1,\Omega_l}
\le\|\Curl\boldsymbol\phi_l\|_{0,\Omega_l}
\le\|\Curl\boldsymbol\phi\|_{0,\Omega}
\qquad\forall l.
\end{equation}
此外, \eqref{eq:sec3-inner-domain-neumann-problem} 等价于
\[
\int_{\Omega_l}\boldsymbol\phi_l\cdot\Grad\mu\dd x=0
\qquad\forall\mu\in H^1(\Omega_l).
\]
所以容易推出
\begin{equation}\label{eq:sec3-inner-domain-l2-bound}
\|\boldsymbol\phi_l\|_{0,\Omega_l}
\le\|\boldsymbol\phi\|_{0,\Omega},
\end{equation}
\begin{equation}\label{eq:sec3-inner-domain-gradient-bound}
\|\Grad\chi_l\|_{0,\Omega_l}
\le\|\boldsymbol\phi\|_{0,\Omega},
\end{equation}
\begin{equation}\label{eq:sec3-inner-domain-gradient-test-bound}
\left|\int_{\Omega_l}\Grad\chi_l\cdot\Grad\mu\dd x\right|
\le\|\boldsymbol\phi\|_{0,\Omega-\Omega_l}
\|\Grad\mu\|_{0,\Omega-\Omega_l}
\quad\forall\mu\in H^1(\Omega).
\end{equation}
把 $\Grad\chi_l$, $\boldsymbol\phi_l$ 和
$\partial\boldsymbol\phi_l/\partial x_i$ 在 $\Omega_l$ 外作零延拓, 并照例用波浪号表示延拓后的函数.
由 \eqref{eq:sec3-inner-domain-gradient-bound} 可知
\[
\widetilde{\Grad\chi_l}\rightharpoonup\mathbf w
\quad\text{weakly in }L^2(\Omega)^3,
\qquad
\Curl\mathbf w=\mathbf0.
\]
事实上, 如果 $\boldsymbol\lambda\in\D(\Omega)^3$, 则存在整数 $k$, 使得
$\boldsymbol\lambda\in\D(\Omega_l)^3$ for all $l\ge k$. 此时

\MathBookSourcePage{70}
\[
\int_\Omega\mathbf w\cdot\Curl\boldsymbol\lambda\dd x
=\lim_{l\to\infty}\int_\Omega
\Grad\chi_l\cdot\Curl\boldsymbol\lambda\dd x,
\]
而结果由下式得到:
\[
\int_{\Omega_l}\Grad\chi_l\cdot\Curl\boldsymbol\lambda\dd x=0
\qquad\forall l\ge k.
\]
由于 $\Omega$ 单连通, 这意味着存在 $\chi\in H^1(\Omega)$, 使得
\[
\mathbf w=\Grad\chi.
\]
此外, \eqref{eq:sec3-inner-domain-gradient-test-bound} 蕴含
$\Grad\chi=\mathbf0$. 因而
\[
\widetilde{\boldsymbol\phi_l}\rightharpoonup\boldsymbol\phi
\quad\text{weakly in }L^2(\Omega)^3.
\]
结合 \eqref{eq:sec3-inner-domain-h1-bound}, 这蕴含
\[
\partial\widetilde{\boldsymbol\phi_l}/\partial x_i
\rightharpoonup\partial\boldsymbol\phi/\partial x_i
\quad\text{weakly in }L^2(\Omega)^3,
\quad1\le i\le3.
\]
所以 $\boldsymbol\phi\in H^1(\Omega)^3$, 并由
\eqref{eq:sec3-inner-domain-h1-bound} 得
\[
|\boldsymbol\phi|_{1,\Omega}
\le\|\Curl\boldsymbol\phi\|_{0,\Omega}.
\]
\end{Proof}

由这个 Theorem 和 Lemma~\ref{lem:sec3-u-embedding-neumann-criterion} 立即得到:
一方面, 如果 $\Omega$ 是有界凸多面体, 问题
\eqref{eq:sec3-u-embedding-neumann-problem} 的解 $\chi$ 自动属于 $H^2(\Omega)$;
另一方面, 如果 $\Omega$ 是非凸多面体, 则有些情形下
\eqref{eq:sec3-u-embedding-neumann-problem} 的解不属于 $H^2(\Omega)$,
从而 $U$ 不包含于 $H^1(\Omega)^3$.

作为 Theorem~\ref{thm:sec3-u-h1-smooth-boundary} 的一个有意义的推论,
有 Foias \& Temam [27] 建立的如下 Corollary.

\setcounter{statement}{6}
\begin{Corollary}\label{cor:sec3-hm-div-curl-characterization}
设 $\Omega$ 是 $\mathbb R^3$ 中的有界开区域, 边界 $\Gamma$ 属于
$\mathcal C^{m,1}$ 类($m\ge1$). 则
\begin{equation}\label{eq:sec3-hm-div-curl-characterization}
\begin{aligned}
H^m(\Omega)^3=\{\mathbf v\in L^2(\Omega)^3;\ &\Div\mathbf v\in H^{m-1}(\Omega),\\
&\Curl\mathbf v\in H^{m-1}(\Omega)^3,\\
&\mathbf v\cdot\mathbf n\in H^{m-1/2}(\Gamma)\},
\end{aligned}
\end{equation}
\begin{equation}\label{eq:sec3-hm-div-curl-estimate}
\|\mathbf v\|_{m,\Omega}\le C\{\|\mathbf v\|_{0,\Omega}
+\|\Div\mathbf v\|_{m-1,\Omega}
+\|\Curl\mathbf v\|_{m-1,\Omega}
+\|\mathbf v\cdot\mathbf n\|_{m-1/2,\Gamma}\}.
\end{equation}
\end{Corollary}
